\documentclass[11pt,a4paper]{article}
\usepackage[T1]{fontenc}
\usepackage[utf8]{inputenc}
\usepackage{lmodern}
\usepackage[indonesian]{babel}
\usepackage[a4paper,margin=2.2cm]{geometry}
\usepackage{microtype}
\usepackage{amsmath,amssymb}
\usepackage{booktabs,tabularx,array}
\usepackage{enumitem}
\usepackage[most]{tcolorbox}
\usepackage{xcolor}
\usepackage{fancyhdr}
\usepackage{needspace}
\usepackage[hidelinks]{hyperref}
\setlength{\parindent}{0pt}
\setlength{\parskip}{5pt}
\setlist{itemsep=2pt,topsep=3pt}
\renewcommand{\arraystretch}{1.25}
\definecolor{biru}{HTML}{1F4E79}
\definecolor{hijau}{HTML}{2E7D32}
\definecolor{merah}{HTML}{B71C1C}
\definecolor{kuning}{HTML}{8D6E00}
\newtcolorbox{intinya}{colback=biru!6,colframe=biru,title=Intinya,fonttitle=\bfseries,boxrule=0.6pt,arc=2pt}
\newtcolorbox{hafal}{colback=hijau!6,colframe=hijau,title=Wajib hafal,fonttitle=\bfseries,boxrule=0.6pt,arc=2pt}
\newtcolorbox{awas}{colback=merah!5,colframe=merah,title=Jangan salah ucap,fonttitle=\bfseries,boxrule=0.6pt,arc=2pt}
\newtcolorbox{contoh}[1][]{colback=kuning!5,colframe=kuning,title={Contoh hitung #1},fonttitle=\bfseries,boxrule=0.6pt,arc=2pt}
\newcounter{tanya}
\newenvironment{tanya}[1]{\refstepcounter{tanya}\par\medskip\needspace{4\baselineskip}\noindent\textbf{\color{biru}T\thetanya.\ #1}\par\noindent\textbf{Jawab:}\ }{\par}
\pagestyle{fancy}\fancyhf{}
\fancyhead[L]{\small Bahan Belajar Proposal --- Abd. Rouf (223140118)}
\fancyhead[R]{\small\thepage}
\renewcommand{\headrulewidth}{0.4pt}
\newcommand{\judul}[2]{\begin{center}{\LARGE\bfseries\color{biru}#1}\par\vspace{4pt}{\large #2}\par\vspace{2pt}{\small Pegangan pribadi untuk bimbingan dan ujian proposal. Isi mengikuti \textit{Proposal\_Skripsi\_ABD\_ROUF.docx}.}\end{center}\vspace{-4pt}\hrule\vspace{6pt}}
